%!TEX root = ../nauchka.tex

\section{Реализация}

\begin{frame}
\frametitle{Архитектура auto\_LiRPA}

\begin{block}{auto\_LiRPA}
Библиотека для автоматического вычисления границ (Linear Relaxation based Perturbation Analysis)
\end{block}

\vspace{0.3cm}

\textbf{Ключевые компоненты:}
\begin{itemize}
    \item \texttt{BoundedModule} --- обёртка над PyTorch/ONNX моделью
    \item \texttt{Bound}-операторы для каждого типа слоя
    \item Методы \texttt{bound\_backward} (CROWN) и \texttt{IBP\_general} (IBP)
\end{itemize}

\end{frame}

\begin{frame}
\frametitle{Реализация TP: операторы и шардирование}

\textbf{Компоненты TP ($\sim$500 строк):}

\begin{enumerate}
    \item \texttt{BoundLinearTP\_Col} / \texttt{BoundLinearTP\_Row}
    \begin{itemize}
        \item Наследуют от \texttt{BoundLinear}, добавляют \texttt{AllReduce} в backward
        \item ONNX-интеграция через \texttt{torch.autograd.Function}
    \end{itemize}

    \item \texttt{tp\_shard\_bounded\_module}
    \begin{itemize}
        \item Автоматическое шардирование произвольного \texttt{BoundedModule}
        \item Замена чередующихся пар \texttt{BoundLinear} на Col/Row
        \item Зонирование подграфов для промежуточных bounds
    \end{itemize}
\end{enumerate}

\end{frame}

\begin{frame}
\frametitle{Реализация FSDP: послойный AllGather}

\textbf{Компоненты FSDP ($\sim$100 строк):}

\begin{enumerate}
    \item \texttt{fsdp\_shard\_bounded\_module}
    \begin{itemize}
        \item Шардирование \texttt{BoundParams} (веса): $W \to W^{(r)} \in \mathbb{R}^{(k/P) \times n}$
    \end{itemize}

    \item Прозрачный \texttt{AllGather} в \texttt{BoundParams.forward()}
    \begin{itemize}
        \item Полная матрица собирается и возвращается как результат forward
    \end{itemize}

    \item Хуки \texttt{fsdp\_gather\_node} / \texttt{fsdp\_free\_node}
    \begin{itemize}
        \item Интегрированы в CROWN backward и IBP forward
        \item В каждый момент в памяти $\leq 1$ полной матрицы (помимо шардов)
    \end{itemize}
\end{enumerate}

\vspace{0.3cm}

\begin{exampleblock}{Ключевое отличие от TP}
Не создаются новые операторы, не изменяется вычислительный граф, не вводятся шардированные зоны
\end{exampleblock}

\end{frame}

\begin{frame}
\frametitle{Статус реализации}

\begin{block}{Tensor Parallelism}
\begin{itemize}
    \item[\checkmark] Операторы \texttt{BoundLinearTP\_Col/Row} с ONNX-интеграцией
    \item[\checkmark] Автоматическое шардирование (PyTorch и ONNX модели)
    \item[\checkmark] Зонирование подграфов, IBP fallback
    \item[\checkmark] Эксперименты: память, корректность, VNN-COMP
\end{itemize}
\end{block}

\vspace{0.2cm}

\begin{block}{FSDP}
\begin{itemize}
    \item[\checkmark] Шардирование весов с послойным AllGather/Free
    \item[\checkmark] Поддержка IBP и CROWN
    \item[\checkmark] Побитово идентичные bounds для всех моделей
    \item[\checkmark] Эксперименты: корректность, память
\end{itemize}
\end{block}

\end{frame}
